\chapter{AI 参与工程作业的证据链}

AI 使用声明不应只问“用了没有”，而应回答“在哪一步用了、输入了什么、怎样核验、谁承担结论”。建议学生随作业提交以下六项。

\begin{enumerate}
  \item \textbf{任务定义：}问题、约束、预期输出和评价标准；
  \item \textbf{来源：}教材、规范、数据与参考模型的版本；
  \item \textbf{AI 记录：}关键提示、工具名称、日期和必要的输出节选；
  \item \textbf{核验：}独立计算、极限情况、量纲、对照案例或第二工具复核；
  \item \textbf{修改：}AI 初稿与最终稿之间最重要的三处变化；
  \item \textbf{责任声明：}仍然存在的不确定性，以及最终判断由谁作出。
\end{enumerate}

\section*{四级量规}

\begin{description}[style=nextline,leftmargin=2.2cm]
  \item[不足] 只有最终成品；来源、AI 参与和核验过程不可追溯。
  \item[基本] 声明工具并保留部分过程，但核验主要停留在“能运行”。
  \item[可靠] 关键假设、单位、边界条件和结果均有独立检查；能指出 AI 错误。
  \item[成熟] 比较多种方案，解释取舍与不确定性，并能在追问中重建判断过程。
\end{description}

量规评的是证据质量，不是提示词长度。提示词写得很长，不代表工程判断可靠。
